This appendix records the pre-registration documents governing the tests
in \S\ref{sec:methods}, together with the commit ordering that makes them
checkable. In every case the decision rules were committed before the
compute they govern was dispatched, and in the DARE-TIES case before the
method being tested had been implemented at all.

\paragraph{Why the ordering is the evidence.} A pre-registration that can
be edited after the fact is worth nothing. The claim we are making is not
that we intended to be honest, but that the repository history makes
dishonesty detectable: each rules commit is an ancestor of every result
commit it governs, and the rules files are not modified after their
initial commit. A reader who does not trust the narrative can verify the
ordering directly.

Table~\ref{tab:commits} gives the hashes, so that checking the ordering does
not require reconstructing which commit is which. Each row reads left to
right in the order the commits were made, and every rules commit is an
ancestor of the result commit on its row.

\paragraph{Two hash spaces, and only one of them resolves.} The
pre-registration documents reproduced in \S\ref{subsec:prereg-docs} cite
commits too, recording by hash that a document was committed before the
generator it governs. Those are the hashes our working repository had when
each document was written, and they will \emph{not} resolve in the published
repository: filtering rewrites every commit, so the hash a document was
written with no longer names anything. We have left them as written rather
than editing the documents, because editing a pre-registration after the fact
is the one thing these documents exist to make detectable, and a reader can
see they are unedited precisely because their internal references are stale.

\textbf{Only the hashes in Tables~\ref{tab:commits} and~\ref{tab:prereg-docs}
resolve in the published repository}, and those are the ones the verification
asks you to use. The correspondence between the two spaces is recoverable:
commit subjects survive filtering, so a stale reference can be matched by its
message.

\begin{table}[tb]
\centering
\small
\caption{Commit trail for every pre-registered test in this paper. The hashes
are those in the public repository, which differ from our working
repository's because the published history is filtered; the ordering is
unaffected, since filtering does not change topology. Verify a row
with \texttt{git merge-base -{}-is-ancestor <rules> <result>}, which exits zero
exactly when the ordering claimed here holds, and fails on the reverse.}
\label{tab:commits}
\begin{tabular}{llll}
\toprule
test & rules & analyzer or implementation & result \\
\midrule
Replication, step 0 (\S\ref{subsec:q1}) & \texttt{947f0ea} &
  --- & \texttt{1dcadd4} \\
Replication, $n=3$ (\S\ref{subsec:q1}) & \texttt{9b3f57e} &
  --- & \texttt{4328446} \\
DARE with TIES (\S\ref{subsec:dare}) & \texttt{5b58059} &
  \texttt{b9b1678}, \texttt{10bd1e4} & \texttt{f33cb02} \\
Conditioning and the ridge (\S\ref{subsec:ridge}) & \texttt{dc0a1f1} &
  \texttt{ac4499f}, \texttt{25d4d47} & \texttt{efa1d2d} \\
Rate exponent (\S\ref{subsec:w3}) & \texttt{dc0a1f1} &
  \texttt{99fd95f} (corrected) & \texttt{c0a1cb3} \\
\addlinespace
Solver replication (\S\ref{subsec:regmean}) & \texttt{876fcb6} &
  \texttt{2adbb79}, \texttt{1733865} & \texttt{57a8f89} \\
Untruncated and gate (\S\ref{subsec:untruncated}) & \texttt{876fcb6} &
  \texttt{e4043b4} & \texttt{63bdf0c} \\
Repaired KnOTS (\S\ref{sec:discussion}) & \texttt{876fcb6} &
  \texttt{45493e4} & \texttt{9f32ed0} \\
Merge matrix at $T=3$ (\S\ref{subsec:q1}) & \texttt{876fcb6} &
  \texttt{4c95967} & \texttt{908a486} \\
\addlinespace
Margin-aware control (\S\ref{subsec:q1}) & \texttt{876fcb6} &
  \texttt{d070a1a} & \texttt{0a4afaf} \\
Shared arm at $n=3$ (\S\ref{subsec:shared-arm}) & \texttt{3b9db2f} &
  \texttt{7c25aa1} & \texttt{0a4afaf} \\
Public-cohort prevalence (\S\ref{subsec:prevalence}) & \texttt{335a4f6} &
  \texttt{0b6eb40} & \texttt{0a4afaf} \\
Downstream accuracy (\S\ref{subsec:downstream-cond}) & \texttt{a56b31e} &
  \texttt{321a808} & \texttt{0a4afaf} \\
Heuristics null in accuracy (\S\ref{subsec:ridge}) & \texttt{91d4a4b} &
  \texttt{4b5a56c} & \texttt{9b77f19} \\
Drift across training (App.~\ref{subsec:drift-result}) & \texttt{8544b36} &
  \texttt{0c8770b} & \texttt{5e01813} \\
\bottomrule
\end{tabular}
\end{table}

Three rows deserve a comment, and the first is the one a sceptical reader
should check first.

The replication appears as two rows because its amendment does \emph{not}
precede the step-0 result: \texttt{9b3f57e} was committed after
\texttt{1dcadd4} was read. That is the blindness limitation the amendment
states about itself, made checkable rather than asserted, and it is why the
amendment is paired with the $n=3$ verdict it actually governs. Running the
ancestry check on \texttt{9b3f57e} against \texttt{1dcadd4} fails, and should.

The DARE-TIES analyzer appears twice because the first version printed, under a
mixed verdict, the consequence text belonging to a different verdict;
\texttt{10bd1e4} is that repair, and it predates the verdict commit. The
rate-exponent row names a \emph{corrected} analyzer, and that correction was
made after the first analysis had been read, for the reason given in
\S\ref{subsec:rate-prereg} below.

\subsection{Replication of the merge matrix on independent cohorts}

Rules committed before any independent-cohort result was read. An
amendment fixing how the three cohorts are combined was committed before
\texttt{indep2} and \texttt{indep3} were read, and states its own
blindness limitation in its first paragraph: \texttt{indep1} had already
been read when the aggregation rule was chosen, so one of three inputs
was visible. We flag this rather than claim full blindness.

The operative rules, summarised here and reproduced in full in the
public repository:
\begin{itemize}
\item Unit of analysis is the cohort, $n = 3$.
\item Tie threshold $0.005$ nats, fixed in the parent document and not
  changed.
\item Noise gate is one-directional: with
  $\mathrm{SE} = \mathrm{sd}/\sqrt{3}$, a win or loss is downgraded to a
  tie unless $|\mathrm{mean}(d)| > 2\times\mathrm{SE}$. The gate may only
  downgrade, never promote.
\item Primary comparison uses the champion baseline selected on the mean
  across cohorts, with two pre-specified robustness lines: champion
  re-selected within each cohort, and the original single-cohort rule
  applied three times with a majority vote.
\item A control question asks how many base models have a non-unanimous
  top-1 method across the three cohorts. Two or more withdraws the
  initialisation-confounding claim.
\end{itemize}

The control returned two of four, and the claim was withdrawn
(\S\ref{subsec:withdrawn}).

\subsection{DARE composed with TIES}

Rules committed before the \texttt{dare\_ties} method existed in the
codebase, which is the strongest ordering guarantee available: the
implementation commit is a descendant of the rules commit, so the rules
cannot have been written to fit an implementation that had not been
written.

The operative rules, summarised here and reproduced in full in the
public repository:
\begin{itemize}
\item $36$ cells: four base models $\times$ three cohorts $\times$ DARE
  density in $\{0.5, 0.2, 0.1\}$, with the TIES trim density pinned at
  $0.2$ in both arms so the only difference from the plain TIES arm is
  the DARE mask.
\item The primary verdict is read at DARE density $0.2$ only, named in
  advance. If another density flatters DARE, that is a secondary
  observation and is not promoted.
\item Per base: helps if mean $d > 0.005$ and the gate passes, hurts if
  mean $d < -0.005$ and the gate passes, neutral otherwise. Overall:
  helps if it helps on $\geq 2$ and hurts on none; hurts if the mirror
  image; neutral if neutral on $\geq 3$; \emph{mixed} otherwise, and then
  neither direction is claimed.
\item The TIES reference cells are the pre-existing ones and are not
  re-run, re-seeded, or re-selected.
\item A smoke test on one real cell is mandatory before the remaining
  thirty-five are dispatched.
\item All three densities are reported, including any that flatter DARE
  and any that do not.
\end{itemize}

The outcome was mixed: helps on one base, hurts on three
(\S\ref{subsec:dare}).

\paragraph{Two disclosures.} First, the mandated smoke cell sits at the
primary density and is therefore one of the twelve cells deciding the
primary verdict, so blindness for that test is eleven of twelve rather
than twelve of twelve. The smoke inspection checked configuration
equality between the arms and that the mask was not a no-op; it did not
examine the direction of the effect. Second, our analysis script
initially printed, under a mixed verdict, the consequence text belonging
to the hurts and neutral verdicts, which asserted that the negative
result generalises. The pre-registration says of a mixed verdict only
that neither direction is claimed. The verdict computation was correct
against the rule; the reporting text was not, and it was corrected before
any of it reached this paper.

\subsection{Rate exponent}\label{subsec:rate-prereg}

The band $[-2.4, -1.6]$ for the log-rate slope, the fit window
$b \in \{2,3,4,8\}$, and the requirement that the band hold on at least
three of four base models were fixed before the sweep ran.

A second document was committed later, when an external review argued that
the flat curve we were then reporting could be an artifact of the
post-merge rank truncation. It fixed three branches in advance: the
falsification is \emph{void} if the slope enters the band on at least
three of four bases, it \emph{stands} if $|{\rm slope}| < 0.5$ on at least
three of four, and it is \emph{unresolved} otherwise, ``reported as
unresolved, not resolved in whichever direction is more convenient''. That
document also recorded, before any new cell ran, what to do about an
anomaly already visible in the output: on two bases the excess at finite
rate sat below the $b = \infty$ value, which the model forbids, and the
rules said in advance that if it persisted it was to be treated as a
harness problem and reported as one.

The experiment those rules were written for was never needed. Its premise
was that the sweep had been truncated, and checking the configurations
showed it had not been: the sweep already ran at the untruncated
realisation, so there was nothing to repeat. The anomaly clause was the
part that mattered, and following it led to the mismatched asymptote and
the refit of \S\ref{subsec:w3}. We applied the three branches to the
refitted slopes as written, obtaining one base in band and none flat,
hence \textsc{unresolved}.

Two things about this ordering deserve stating plainly. The decision rule
is unchanged and predates the numbers it judges. But the rule was
registered against a different intervention from the one that ended up
producing those numbers, so what carried over is the standard of evidence
rather than a full pre-specification of the analysis, and the correction
itself was made after the first analysis had been read. We think a
forbidden value obliges investigation rather than permitting it, which is
why the clause was written in advance, and the commit ordering shows both
the clause and what followed it.

\subsection{The pre-registration documents}\label{subsec:prereg-docs}

Twelve documents govern the tests in this paper. Table~\ref{tab:prereg-docs}
lists them with the commit that introduced each, and the rules are
summarised below. \textbf{The documents themselves are in the public
repository, complete and unedited}, at the commits named here, which is the
copy that cannot have been written afterwards. Earlier versions of this paper
reproduced them all verbatim in the appendix; that is what the repository is
for, and printing them twice made the paper longer without making anything
more checkable.

Each is the file as committed and none was modified after its commit, which
is checkable with \texttt{git log -{}-follow} on the path. One character is
altered across the twelve: the first document records the wall-clock time at
which a set of cells completed, with a timezone abbreviation that would
narrow our geography under double-blind review, and it appears as
\texttt{[tz]}. That is the only alteration.

Two conventions are shared by every document and are not repeated in the
summaries: a \textbf{tie threshold of $0.005$ nats}, fixed in the
2026-08-03 parent, and a \textbf{one-directional noise gate} at
$2\times\mathrm{SE}$ over cohorts, which may downgrade a win or a loss to a
tie but may never promote a tie to either. Verdicts are reported in
whichever branch the rule lands in.

\paragraph{The rules, one paragraph each.}
\begin{itemize}\itemsep2pt
\item \textbf{\texttt{prereg\_a1\_matrix}} (\texttt{947f0ea}) fixes the
replication: does the encoder's win over the champion baseline survive the
move to independently initialised cohorts, and do method rankings change
between cohorts? It states its own statistical ceiling up front, that one
cohort per arm cannot support a strong verdict regardless of what the rules
return. Outcome: the win does not survive (\S\ref{subsec:q1}).

\item \textbf{\texttt{prereg\_a1\_matrix\_amendment}} (\texttt{9b3f57e})
extends it to $n = 3$ cohorts and fixes the aggregation. It declares its
own blindness status plainly: \texttt{indep1} had been read when the rule
was written, so the rule was chosen to be the obvious one rather than
selected from a menu. It also retires the parent's $0.013$ nat provisional
band.

\item \textbf{\texttt{prereg\_dare\_ties}} (\texttt{5b58059}) asks whether
our DARE-with-task-arithmetic result extends to DARE with TIES, primary at
density $0.2$ only, with the density sweep explicitly secondary so that no
density can be chosen after the fact. Outcome: it does not extend
(\S\ref{subsec:dare}).

\item \textbf{\texttt{prereg\_conditioning}} (\texttt{dc0a1f1}) governs the
ridge sweep and the rate exponent: whether the optimal ridge tracks the
conditioning, and whether the flat rate curve is a truncation artifact. It
fixes the exponent band at $[-2.4, -1.6]$ on at least three of four bases.
Outcomes: the ridge \emph{gain} tracks conditioning while $\lambda^\star$
does not (\S\ref{subsec:ridge}), and the exponent is unresolved
(Appendix~\ref{app:rate}).

\item \textbf{\texttt{prereg\_tmlr}} (\texttt{876fcb6}) is the revision
campaign: the go/no-go is whether cohort conditioning affects a solver we
did not build, with the two cohorts required to match byte-for-byte on
tasks, evaluation seed, sequence length and VRAM gate. It also governs the
KnOTS repair and the $T = 3$ arm. Outcome: both predictions hold on four of
four bases (\S\ref{subsec:regmean}).

\item \textbf{\texttt{prereg\_tmlr\_amendment}} (\texttt{2adbb79}) sets
RegMean's minimally regularised reference to $\lambda = 10^{-6}$ rather
than $0$, because RegMean solves in the full input dimension against a Gram
of rank at most $Tr$, so $\lambda = 0$ is singular rather than
ill-conditioned. It changes a grid point, not a threshold or a rule.

\item \textbf{\texttt{prereg\_shared\_arm}} (\texttt{3b9db2f}) answers a
limitation the draft raised against itself: the ridge sweep's shared arm
was $n = 1$, so its gate was one-sided. This registers the shared arm at
$n = 3$ with a genuinely two-sided gate, and flags in advance that the
published cohort sitting more than $2$ sd from its arm's mean would
require reporting. Outcome: both predictions $4/4$, and the published
cohort sits $0.23$ to $0.75$ sd out (\S\ref{subsec:shared-arm}).

\item \textbf{\texttt{prereg\_prevalence}} (\texttt{335a4f6}) fixes the
sampling frame for the public-cohort audit before any download, in three
parts: the library condition (A), published benchmarks' released cohorts
(B), and the broad sample (C). Outcome: A and C reported
(\S\ref{subsec:prevalence}), B not run.

\item \textbf{\texttt{prereg\_downstream}} (\texttt{a56b31e}) asks whether
the conditioning collapse is visible in accuracy rather than only in NLL,
on GSM8K and HumanEval, with a binomial-SE gate. It fixes in advance what
we would write in each branch, including that a null would move the metric
limitation into \S\ref{sec:intro} and delete the practice recommendations.
Outcome: the ridge recovers the shared arm on four of four bases on both
benchmarks (\S\ref{subsec:downstream-cond}).

\item \textbf{\texttt{prereg\_downstream\_amendment}} (\texttt{5216230})
corrects our own specification error: it registered $n = 500$ evaluation
items, but HumanEval has $164$. The gate adapts to the actual $n$, which
widens it to roughly $11$ points on HumanEval against $6$ on GSM8K, so the
two benchmarks are not equally powered and a null on HumanEval is not read
as equivalent to a null on GSM8K. It makes our own predictions harder to
satisfy, and was written before any cell ran.

\item \textbf{\texttt{prereg\_drift}} (\texttt{8544b36}) discharges E3,
which asked how far $A$ travels from its initialisation across training and
which we had reported as unevaluable because the original cohorts kept no
initial weights and no intermediate checkpoints. It trains a fresh cohort with
both, on one base model in both initialisation arms, and fixes three
predictions: that $A$ ends within half its own norm of where it started, that
the trace is monotone, and that the cohort reproduces the collapse at all.
Branch~3 is the one that would have cost us: if $A$ travelled far and the
subspaces still collapsed, the mechanism of \S\ref{subsec:mechanism} would be
falsified and the abstract would have to say so. Outcome: all three hold
(Appendix~\ref{subsec:drift-result}).

\item \textbf{\texttt{prereg\_heuristics\_downstream}} (\texttt{91d4a4b})
answers an asymmetry we found in our own reporting: the null of
\S\ref{subsec:ridge} is measured only in worst-task NLL excess, the unit
\S\ref{sec:discussion} disclaims, while the positive result beside it has
two units. It registers the same five heuristics on the same two cohorts,
scored on GSM8K and HumanEval, $80$ cells. Two features are worth naming
because both cut against us. The gate is \emph{binomial} rather than the
$2\times\mathrm{SE}$-over-cohorts gate used elsewhere, since this design has
one cohort per arm, which makes it a weaker instrument with a minimum
detectable effect of about $9$ accuracy points on GSM8K and $16$ on
HumanEval; the document fixes the wording any null must be reported in so
that it cannot later be read as ``no effect''. And it re-runs \emph{both}
arms rather than reusing the $40$ shared-arm cells that already exist,
because those were scored before the scorer repairs of
\S\ref{subsec:downstream-cond} and a paired difference taken across two
scorer versions would carry a known, method-dependent defect inside it. All
four branches of the decision rule are written out in advance, and two of
them remove or narrow contribution~3.
\end{itemize}

\begin{table}[htb]
\centering
\small
\caption{The twelve pre-registration documents and the commit that introduced
each. Every one was added in a single commit and never modified afterwards,
which is checkable with \texttt{git log -{}-follow} on the path.}
\label{tab:prereg-docs}
\begin{tabular}{lll}
\toprule
document & introduced in & governs \\
\midrule
\texttt{prereg\_a1\_matrix} & \texttt{947f0ea} & replication, step 0 \\
\texttt{prereg\_a1\_matrix\_amendment} & \texttt{9b3f57e} & replication, $n=3$ \\
\texttt{prereg\_dare\_ties} & \texttt{5b58059} & DARE with TIES \\
\texttt{prereg\_conditioning} & \texttt{dc0a1f1} & ridge sweep, rate exponent \\
\texttt{prereg\_tmlr} & \texttt{876fcb6} & solver replication, gate, KnOTS, $T=3$ \\
\texttt{prereg\_tmlr\_amendment} & \texttt{2adbb79} & RegMean's minimal $\lambda$ \\
\texttt{prereg\_shared\_arm} & \texttt{3b9db2f} & the shared arm at $n=3$ \\
\texttt{prereg\_prevalence} & \texttt{335a4f6} & public-cohort prevalence \\
\texttt{prereg\_downstream} & \texttt{a56b31e} & downstream accuracy \\
\texttt{prereg\_downstream\_amendment} & \texttt{5216230} & HumanEval's 164 items \\
\texttt{prereg\_heuristics\_downstream} & \texttt{91d4a4b} & the heuristics null in accuracy \\
\texttt{prereg\_drift} & \texttt{8544b36} & how far $A$ travels (E3) \\
\bottomrule
\end{tabular}
\end{table}

